\section{Case Study: The Evolution of Narrative and Propaganda Detection}
The field of automated narrative and propaganda detection serves as an ideal case study, illustrating the practical application and evolution of the very techniques discussed in this survey. This domain has acted as a methodological crucible, progressing from standard classification to sophisticated, multilingual, and hierarchical analysis, with the premier NLP benchmarking venue, SemEval, providing a clear chronological roadmap of the state-of-the-art.

\subsection{From Flat Classification to Structural Priors}
The initial wave of research was benchmarked by SemEval-2020 Task 11, which framed propaganda detection as a flat multi-class classification problem with 14 techniques~\cite{SemEval2020Task11}. The results unequivocally established the dominance of fine-tuned Transformer-based architectures like RoBERTa. However, top systems rarely used these models off-the-shelf. For identifying propagandistic text spans, a sequence tagging sub-task, the most effective architecture was a RoBERTa model augmented with a Conditional Random Field (CRF) layer, which adds a structural constraint by modeling dependencies between adjacent labels~\cite{chernyavskiy-etal-2020-aschern, jurkiewicz-etal-2020-applicaai}. For the technique classification sub-task, the winning solutions relied on large ensembles of RoBERTa models to improve robustness and accuracy~\cite{jurkiewicz-etal-2020-applicaai}.

Subsequent research sought to improve performance by injecting deeper linguistic knowledge, moving beyond surface-level features. A landmark contribution was the discourse-aware Transformer architecture proposed by Chernyavskiy et al. (2024), which integrates features from Rhetorical Structure Theory (RST)~\cite{chernyavskiy-etal-2024-unleashing}. By providing the model with explicit information about discourse units (EDUs), their salience (Nucleus vs. Satellite), and their relationships, the model gained a stronger inductive bias. Experiments showed that propaganda was most frequently located in the 'Nucleus' EDUs---the core parts of a message. This work provides compelling evidence that incorporating structural linguistic priors is a powerful method for enhancing models for persuasive language analysis, particularly in low-resource settings where such structures cannot be learned from data alone.

\subsection{The Hierarchical and Multilingual Turn in SemEval}
The challenges of multilinguality, data imbalance, and hierarchical classification, which are central to modern HMLC, were addressed progressively in subsequent SemEval shared tasks.

\textbf{SemEval-2023 Task 3} dramatically expanded the scope to nine languages and 23 persuasion techniques~\cite{SemEval2023Task3}. This pivot forced a move from English-centric models to multilingual encoders like XLM-RoBERTa. The winning strategies provided a clear blueprint for low-resource multilingual classification: 1) joint training on data from all languages to learn shared cross-lingual representations, and 2) aggressive data augmentation via machine translation, where examples from high-resource languages are translated to bolster the training sets of low-resource ones~\cite{hromadka-etal-2023-kinitveraai, modzelewski-etal-2023-dshacker}.

\textbf{SemEval-2024 Task 4} marked the formal introduction of hierarchy, reframing persuasion detection as an HMLC problem with 22 techniques organized in a Directed Acyclic Graph (DAG)~\cite{SemEval2024Task4}. This task spurred the development of novel hierarchy-aware architectures, with top systems showcasing three distinct state-of-the-art approaches:
\begin{itemize}
    \item \textbf{Direct HTC Modeling:} Adapting dedicated Hierarchical Text Classification models like HPT, which required pre-processing the label DAG into a strict tree structure~\cite{HierarchyEverywhere2024}.
    \item \textbf{Architectural Priors:} Designing a custom hierarchical classification head on top of an LLM backbone. This head mirrors the label hierarchy with structured layers, forcing predictions to be made in a top-down, consistent manner~\cite{wunderle-etal-2024-otterlyobsessedwithsemantics}. This is a direct implementation of the local, top-down approach discussed in Section~\ref{sec:local_vs_global}.
    \item \textbf{Representation Learning:} Employing hyperbolic embeddings to project text and label representations into a non-Euclidean space whose geometry naturally mirrors a tree-like structure, allowing the model to implicitly learn hierarchical relationships~\cite{chikoti-etal-2024-iitk}.
\end{itemize}

Finally, \textbf{SemEval-2025 Task 10} represents the culmination of this trend, establishing multilingual narrative classification as a formal shared task with a two-level hierarchical taxonomy~\cite{semeval2025task10}. Early systems already demonstrate the application of advanced techniques like retrieval-based methods and hierarchical prompting with LLMs~\cite{younus-qureshi-2025-nlptuducd, singh-etal-2025-gatenlp}. This clear evolutionary trajectory within a major NLP benchmark provides an exceptionally strong justification for research into HMLC for narrative detection, as it confirms that this problem is the logical and timely next step for the research community.